% Рисунок с доски: psi = 0, начальный профиль phi -- треугольник. Он распадается
% на две бегущие волны половинной высоты со скоростями -a и +a.
\begin{figure}[htbp]
\centering
\begin{tikzpicture}
\begin{groupplot}[
  group style={group size=2 by 1, horizontal sep=1.9cm},
  width=6.6cm, height=4.2cm,
  axis lines=middle,
  xlabel={$x$},
  xlabel style={anchor=north west},
  xtick=\empty, ytick=\empty, clip=false,
]
% --- слева: начальный профиль
\nextgroupplot[xmin=-2.6, xmax=2.6, ymin=-0.35, ymax=1.45]
  \addplot[very thick, blue!60!black] coordinates {(-2.4,0) (-0.6,0) (0,1) (0.6,0) (2.4,0)};
  \node[blue!60!black, anchor=south east, font=\footnotesize] at (axis cs:-0.1,0.78) {$\varphi(x)$};
  \node[anchor=south, font=\footnotesize] at (axis cs:0,1.12) {$\psi \equiv 0$};
  \node[anchor=east, font=\footnotesize] at (axis cs:-0.04,0.98) {$1$};

% --- справа: момент времени t > 0
\nextgroupplot[xmin=-2.6, xmax=2.6, ymin=-0.35, ymax=1.45,
               xtick={-1.4,1.4}, xticklabels={$-at$,$at$}]
  \addplot[very thick, blue!60!black]
    coordinates {(-2.4,0) (-2.0,0) (-1.4,0.5) (-0.8,0) (0.8,0) (1.4,0.5) (2.0,0) (2.4,0)};
  \draw[dashed, gray] (axis cs:-1.4,0) -- (axis cs:-1.4,0.5);
  \draw[dashed, gray] (axis cs: 1.4,0) -- (axis cs: 1.4,0.5);
  \draw[-{Latex[length=4pt]}, gray] (axis cs:-1.0,0.75) -- (axis cs:-1.7,0.75);
  \node[anchor=south, font=\footnotesize] at (axis cs:-1.35,0.78) {$a$};
  \draw[-{Latex[length=4pt]}, gray] (axis cs:1.0,0.75) -- (axis cs:1.7,0.75);
  \node[anchor=south, font=\footnotesize] at (axis cs:1.35,0.78) {$a$};
  \node[anchor=east, font=\footnotesize] at (axis cs:-1.44,0.5) {$\tfrac12$};
\end{groupplot}
\end{tikzpicture}
\caption{Случай $\psi \equiv 0$: начальный профиль распадается на две волны половинной
амплитуды, бегущие влево и вправо со скоростью $a$. Их вершины находятся в точках
$x = -at$ и $x = at$.}
\label{fig:wave-split}
\end{figure}
